%% ============================================================
%% fhe/parameter-security.tex
%% ============================================================
%%
%% Canonical Lux FHE LWE-parameter security theorem. Papers cite
%% this; do not redefine.
%% ============================================================

\subsection*{Lux FHE concrete LWE-parameter security}

\begin{theorem}[Concrete LWE security under the Albrecht estimator]
\label{thm:lwe-parameter-security}
Let $\mathsf{ps} \in \{\mathsf{PN10QP27}, \mathsf{PN11QP54},
\mathsf{PN9QP28\_STD128}\}$ denote a Lux FHE production parameter set with
LWE dimension $n = 2^{\log_2 n_{\mathrm{LWE}}(\mathsf{ps})}$, modulus
$q = q_{\mathrm{LWE}}(\mathsf{ps})$, and noise standard deviation
$\sigma_e(\mathsf{ps})$ (the latter resolved by the parameter-set keygen
schedule). Then under the Albrecht-Player-Scott LWE estimator with the
BKZ block-size cost model and the BDGL sieving cost,
\[
  \mathrm{Adv}^{\mathsf{LWE}}_{n, q, \sigma_e}(\mathcal{A})
  \;\leq\;
  2^{-\lambda(\mathsf{ps})}
\]
with $\lambda(\mathsf{ps}) \geq 128$ classical bits and $\lambda(\mathsf{ps})
\geq 128$ quantum bits (under Grover speedup) for every shipped
$\mathsf{ps}$.
\end{theorem}

\begin{remark}[Estimator citation]
\label{rem:albrecht-cite}
The estimator is Albrecht, Player, and Scott, "On the concrete hardness
of Learning with Errors", J.\ Math.\ Cryptology 9(3):169--203, 2015,
together with Albrecht's open-source \texttt{lwe-estimator} updates as of
2026-Q1. Concrete numbers per parameter set are reported in
\texttt{papers/lux-fhe-runtime/07-benchmarks.tex} alongside the
benchmark hardware.
\end{remark}

\begin{remark}[NIST IR 8214C alignment]
\label{rem:nist-8214c-alignment}
The threshold-FHE keygen path (LP-019, LP-076, LP-167 \S Key custody)
tracks the NIST Threshold Cryptography Project's evaluation framing
(NIST IR 8214C). Lux FHE parameter-set lambda targets are chosen to
exceed NIST's category-1 floor under both classical and quantum
adversaries.
\end{remark}

\begin{remark}[No PQ rollback]
\label{rem:no-pq-rollback}
LWE / RLWE security is the post-quantum security ground truth for the
Lux FHE stack; Lux does not weaken parameters for performance. The CPU
fallback (\texttt{luxcpp/crypto/fhe/cpp/backends/cpu/}) uses identical
parameters to the GPU backends. The dual-runtime equivalence theorem
(\texttt{proofs/fhe/dual-runtime-equivalence.tex},
Theorem~\ref{thm:fhe-dual-runtime-equivalence}) ensures that whichever
backend the operator selects, the underlying LWE-hardness assumption
applies unchanged.
\end{remark}
